\subsection{Introduzione}
\begin{rema}
Nel linguaggio Wolfram non ci sono tipi: l'unico tipo è il tipo \textit{funzione}.
Non c'è memoria, ma solo \textit{stati} (come per gli automi).
Si possono aggiungere stati e/o regole (queste ultime all'interno del ``kernel'' di Wolfram).
\end{rema}
\begin{defi}
Una \textit{espressione} è un oggetto della forma
\[
E[E_1, \dots, E_n]
\]
dove, a loro volta, gli $E_i$ sono espressioni.
\end{defi}
\begin{defi}
Data una funzione $f$ e un'espressione $E[E_1,\dots,E_n]$, si può costruire
\[
\wolfram{Map}[f, E[E_1, \dots, E_n]],
\]
che applica $f$ a ciascun argomento $E_i$, mantenendo la testa $E$.
\end{defi}
\begin{rema}
Ad esempio,
\[
\wolfram{Map}[f,\{a_1, \dots, a_n\}]
= \wolfram{Map}[f,\wolfram{List}[a_1, \dots, a_n]]
= \wolfram{List}[f[a_1], \dots, f[a_n]].
\]
\end{rema}
\begin{defi}
Ci sono due modi per introdurre regole nel kernel.
\begin{enumerate}
\item \textit{Assegnazione immediata} (\wolfram{=}): il membro destro viene \emph{valutato} e la regola di sostituzione così ottenuta viene inserita per tutte le espressioni che hanno la stessa forma:
\[
\wolfram{f[x_] = x + 2*a}
\]
\item \textit{Assegnazione ritardata} (\wolfram{:=}): il membro destro \emph{non} viene valutato, ma si memorizza soltanto la regola, che resta scritta così com'è:
\[
\wolfram{f[x_] := x + 2*a}
\]
\end{enumerate}
\end{defi}
\begin{rema}
Nel secondo caso la valutazione avviene solo quando si chiama, ad esempio, \wolfram{f[3]}: in quel momento vengono ricalcolate tutte le operazioni usando l'ultimo valore assegnato anche a tutti i parametri presenti nella regola.
\end{rema}
\begin{rema}
Esempio di codice:
\begin{walfram}
In[1]:=  f[x_] = 2x + a
Out[1]=  2x + a
In[2]:=  f[3]
Out[2]=  6 + a
In[3]:=  a = 2
Out[3]=  2
In[4]:=  f[3]
Out[4]=  8
\end{walfram}
Con \wolfram{=} il membro destro è stato valutato subito (con $a$ ancora indeterminato), mentre $a$ viene sostituito solo alla chiamata successiva.
\end{rema}
\begin{rema}
Per vedere le regole definite su un simbolo $f$ basta digitare
\[
\wolfram{?f}
\]
\end{rema}
\begin{defi}
Con \wolfram|List[1,2,3]| (o \wolfram|{1,2,3}|) si crea l'insieme ordinato $\{1,2,3\}$.
\end{defi}
\begin{defi}
Il comando
\[
\wolfram|Apply[k, {1,2,3}]|
\]
sostituisce la testa \wolfram{List} con $k$, dando come output
\[
k[1,2,3].
\]
\end{defi}
\begin{rema}
\wolfram{Plus} fa la somma degli elementi, quindi
\[
\wolfram|Apply[Plus, {1,2,3}]| \quad\longmapsto\quad 6.
\]
Cambiando \wolfram{Plus} in \wolfram{Times} si ottiene la produttoria:
\[
\wolfram|Apply[Times, {1,2,3}]| \quad\longmapsto\quad 6.
\]
\end{rema}
\begin{rema}
Per eseguire più istruzioni in un certo ordine, le si separa con un ``\wolfram{;}''. Il punto e virgola sopprime anche l'output dell'istruzione che precede.
\end{rema}
\begin{rema}
Esempio di codice:
\begin{walfram}
lista = Range[0, 1000]; Map[h, lista]
\end{walfram}
Il comando \wolfram{Range}
crea un'insieme
ordinato (in questo caso da $0$ a $1000$).
L'output è la lista $\{h[0], h[1], \dots, h[1000]\}$.
\end{rema}
\begin{defi}
Per estrarre elementi da una lista (o più in generale i rami dell'albero di un'espressione) si usa il comando \wolfram{Part}, comunemente abbreviato con le doppie parentesi quadre \wolfram{[[ ]]}. 
Ad esempio, se \wolfram{lista = \{a, b, c\}}, il comando \wolfram{lista[[2]]} restituisce \wolfram{b}. 
È possibile usarlo in modo annidato: \wolfram{matrice[[1, 2]]} estrae il secondo elemento della prima riga. In \wolfram{TM-lib.m}, \wolfram{automa[[1]]} estrae la tabella delle transizioni.
\end{defi}

\begin{defi}
In Wolfram è possibile definire \textit{funzioni pure} (o anonime) senza assegnare loro un nome. Si utilizza il simbolo \wolfram{\#} per indicare l'argomento (o gli argomenti, es. \wolfram{\#1, \#2}) e il simbolo \wolfram{\&} alla fine per indicare la chiusura della funzione. Ad esempio:
\[
\wolfram{Map[\#^2 \&, \{1, 2, 3\}]} \quad\longmapsto\quad \{1, 4, 9\}.
\]
Questo costrutto è estremamente comune per applicare trasformazioni rapide all'interno di un \wolfram{Map}.
\end{defi}

\begin{defi}
Oltre all'assegnazione globale 
(\wolfram{=}), esistono le \textit{regole di trasformazione} locali, create con \wolfram{->} (\wolfram{Rule}). Per applicare una o più regole a un'espressione si usa l'operatore di sostituzione \wolfram{/.} (\wolfram{ReplaceAll}). Ad esempio:
\[
\wolfram{x^2 + y  /.  x -> 3} \quad\longmapsto\quad 9 + y.
\]
Negli automi, le singole transizioni sono modellate proprio come liste di regole del tipo \wolfram{\{stato, simbolo\} -> nuovo\_stato}.
\end{defi}

\begin{rema}
Hai già visto che il comando \wolfram{Apply[f, espressione]} sostituisce la testa dell'espressione con \wolfram{f}. Esiste una forma sintattica abbreviata e molto diffusa per questa operazione, che usa il simbolo \wolfram{@@}. Pertanto:
\[
\wolfram{Apply[Plus, \{1, 2, 3\}]} \quad \text{è equivalente a} \quad \wolfram{Plus @@ \{1, 2, 3\}}.
\]
Nel codice della libreria, ad esempio, si usa \wolfram{Union @@ ...} per estrarre e unire in un unico insieme gli elementi di una lista di liste.
\end{rema}
\begin{defi}
In Wolfram Language, il costrutto \wolfram{Module} viene utilizzato per definire variabili locali. La sintassi generale è:
\[
\wolfram{Module[\{var1, var2, \dots\}, corpo\_della\_funzione]}
\]
Questo ambiente garantisce che le variabili dichiarate nella lista iniziale (ad esempio \wolfram{var1}, \wolfram{var2}) abbiano visibilità solo all'interno del modulo, evitando che vadano a sovrascrivere o interferire con variabili globali omonime definite altrove nel kernel.
\end{defi}

\begin{rema}
Il blocco \wolfram{Module} raggruppa più istruzioni separate da punto e virgola (\wolfram{;}) e restituisce come output il valore dell'ultima espressione valutata. Ad esempio:
\begin{walfram}
calcola[x_] := Module[{y, z},
  y = x^2;
  z = y + 1;
  z * 2
]
\end{walfram}
In questo caso, \wolfram{y} e \wolfram{z} sono variabili temporanee locali. L'output della chiamata \wolfram{calcola[3]} sarà il risultato dell'ultima operazione, ovvero \wolfram{20}.
\end{rema}
\subsection{Entropia}
\begin{defi}
Entropia \dots % da completare
\end{defi}
\begin{defi}
Sia $(\Omega,\mathcal{F},\mathbb{P})$ uno spazio di probabilità e sia $\mathcal{X}$ un insieme finito (l'alfabeto), munito della $\sigma$-algebra $2^{\mathcal{X}}$. Una variabile aleatoria è una funzione misurabile $X\colon \Omega \to \mathcal{X}$. La sua legge è la misura di probabilità $\mathbb{P}_X = \mathbb{P}\circ X^{-1}$ su $\mathcal{X}$, e si pone
\[
p(x) = \mathbb{P}_X(\{x\}) = \mathbb{P}(X = x), \qquad x \in \mathcal{X}.
\]
\end{defi}
\begin{defi}
Il \textit{contenuto informativo} (grezzo) di un esito $x\in\mathcal{X}$ con $p(x)>0$ è
\[
h(x) = \log \frac{1}{p(x)} = -\log p(x).
\]
Con il logaritmo in base $2$ l'unità di misura è il bit. Più un esito è improbabile, maggiore è la sua informazione: $h(x)=0$ se e solo se $p(x)=1$.
\end{defi}
\begin{defi}
Date due variabili aleatorie $X\colon\Omega\to\mathcal{X}$ e $Y\colon\Omega\to\mathcal{Y}$, la loro legge congiunta è $p(x,y)=\mathbb{P}(X=x,\,Y=y)$. Il contenuto informativo della coppia è
\[
h(x,y) = \log \frac{1}{p(x,y)}.
\]
\end{defi}
\begin{defi}
Le variabili $X$ e $Y$ sono \textit{indipendenti} se $\sigma(X)$ e $\sigma(Y)$ sono indipendenti, cioè $\mathbb{P}(X\in A,\, Y\in B)=\mathbb{P}(X\in A)\,\mathbb{P}(Y\in B)$ per ogni $A\subseteq\mathcal{X}$ e $B\subseteq\mathcal{Y}$. Poiché gli alfabeti sono finiti, ciò equivale a
\[
p(x,y) = p(x)\,p(y) \qquad \forall\, x\in\mathcal{X},\ y\in\mathcal{Y}.
\]
\end{defi}
\begin{rema}
Se $X$ e $Y$ sono indipendenti, allora l'informazione della coppia è la somma delle informazioni:
\[
h(x,y) = \log\frac{1}{p(x)\,p(y)} = \log\frac{1}{p(x)} + \log\frac{1}{p(y)} = h(x) + h(y).
\]
Questa è la ragione per cui si usa il logaritmo: trasforma il prodotto delle probabilità (indipendenza) in una somma di informazioni.
\end{rema}
\begin{defi}
L'\textit{entropia} di $X$ è il valore atteso del contenuto informativo:
\[
H(X) = \mathbb{E}\big[h(X)\big] = \int_{\Omega} h\big(X(\omega)\big)\, d\mathbb{P}(\omega).
\]
Per il teorema di cambio di variabili rispetto alla legge $\mathbb{P}_X$ si ottiene
\[
H(X) = \int_{\mathcal{X}} h(x)\, d\mathbb{P}_X(x) = \sum_{x\in\mathcal{X}} p(x)\, h(x) = \sum_{x\in\mathcal{X}} p(x)\log\frac{1}{p(x)},
\]
con la convenzione $0\log\frac{1}{0}=0$, che permette di sommare anche sugli $x$ con $p(x)=0$.
\end{defi}
\begin{rema}
L'entropia misura l'informazione media, o equivalentemente l'incertezza, associata a $X$. Dipende solo dalla legge $\mathbb{P}_X$ e non dai valori assunti da $X$. Vale sempre $H(X)\ge 0$, poiché $h(x)\ge 0$ quando $p(x)>0$. Inoltre, essendo $\mathcal{X}$ finito, $h(X)$ è una funzione semplice e quindi integrabile, dunque $H(X)<\infty$.
\end{rema}
\begin{rema}
Se $X$ e $Y$ sono indipendenti, $H(X,Y)=H(X)+H(Y)$. Infatti, per la linearità del valore atteso e l'additività di $h$,
\[
H(X,Y) = \mathbb{E}\big[h(X,Y)\big] = \mathbb{E}\big[h(X)+h(Y)\big] = \mathbb{E}\big[h(X)\big] + \mathbb{E}\big[h(Y)\big] = H(X) + H(Y).
\]
\end{rema}
\begin{rema}
Esempio di calcolo dell'entropia (in bit) di una distribuzione $p=(\tfrac12,\tfrac14,\tfrac14)$:
\begin{walfram}
p = {1/2, 1/4, 1/4};
Total[p * Log[2, 1/p]]
\end{walfram}
L'output è $\tfrac{3}{2}$, infatti $\tfrac12\cdot 1+\tfrac14\cdot 2+\tfrac14\cdot 2=\tfrac32$.
\end{rema}
\begin{defi}
Appendice: fondamenti di probabilità. Uno \textit{spazio di probabilità} è una terna $(\Omega,\mathcal{F},\mathbb{P})$ dove:
\begin{itemize}
\item $\mathcal{F}\subseteq 2^{\Omega}$ è una $\sigma$-algebra, cioè $\Omega\in\mathcal{F}$, è chiusa per complementare e per unioni numerabili;
\item $\mathbb{P}\colon\mathcal{F}\to[0,1]$ è una misura con $\mathbb{P}(\Omega)=1$, cioè è $\sigma$-additiva sugli elementi disgiunti di $\mathcal{F}$.
\end{itemize}
Una funzione $X\colon(\Omega,\mathcal{F})\to(E,\mathcal{E})$ è \textit{misurabile} se $X^{-1}(A)\in\mathcal{F}$ per ogni $A\in\mathcal{E}$. La \textit{legge} di $X$ è $\mathbb{P}_X(A)=\mathbb{P}(X^{-1}(A))$ per $A\in\mathcal{E}$.
\end{defi}
\begin{defi}
Sia $g\colon\Omega\to\mathbb{R}$ misurabile. Il \textit{valore atteso} di $g$ è l'integrale di Lebesgue
\[
\mathbb{E}[g] = \int_{\Omega} g\, d\mathbb{P},
\]
definito prima per le funzioni semplici $g=\sum_{i} c_i \mathbbm{1}_{A_i}$ come $\sum_i c_i\,\mathbb{P}(A_i)$, poi per $g\ge0$ come estremo superiore sulle funzioni semplici $0\le s\le g$, e infine per $g$ integrabile come $\mathbb{E}[g^+]-\mathbb{E}[g^-]$.
\end{defi}
\begin{rema}
Vale il \textit{cambio di variabili}: se $X\colon\Omega\to E$ è misurabile e $f\colon E\to\mathbb{R}$ è misurabile e integrabile rispetto a $\mathbb{P}_X$, allora
\[
\mathbb{E}\big[f(X)\big] = \int_{\Omega} f\circ X\, d\mathbb{P} = \int_{E} f\, d\mathbb{P}_X.
\]
Se $E$ è finito o numerabile, l'ultimo integrale si riduce alla somma $\sum_{x\in E} f(x)\,\mathbb{P}_X(\{x\})$. Questa è la formula usata per scrivere $H(X)=\sum_x p(x)h(x)$.
\end{rema}
\begin{rema}
Calcoliamo \wolfram|Integrate[x^2 Sin[x], x]| e vediamo il risultato:
\[
\int x^2 \sin x \, dx = (2 - x^2)\cos x + 2x \sin x + C.
\]
\end{rema}
\subsection{Libreria TM-lib}
\begin{defi}
Per implementare gli automi in Wolfram si usa il pacchetto \wolfram{TM-lib.ma}. Si procede in tre passi.
\begin{enumerate}
\item Assegnare la directory corrente, cioè la cartella che contiene il pacchetto: \wolfram|SetDirectory["/percorso/della/cartella"]|
\item Caricare la libreria: \wolfram|Get["TM-lib.ma"]|
\item Usare le funzioni implementate dalla libreria (elencate sotto).
\end{enumerate}
\end{defi}
\begin{rema}
Esempio di codice per i primi due passi:
\begin{walfram}
SetDirectory["/percorso/della/cartella"]
Get["TM-lib.ma"]
\end{walfram}
\end{rema}
\begin{defi}
La libreria implementa le seguenti funzioni.
\begin{itemize}
\item \wolfram|ShowConfiguration[c_]| visualizza la configurazione (di una macchina di Turing).
\item \wolfram|ManipulateComputation[mu_, init_]| visualizza la computazione di una macchina di Turing.
\item \wolfram|InitAutomaton[automaton_, w_]| inizializza l'automa, cioè restituisce la configurazione iniziale associata alla parola $w$.
\item \wolfram|DFA2TM[automaton_]| trasforma un automa nella macchina di Turing equivalente.
\item \wolfram|States[automaton_]| restituisce la lista degli stati dell'automa.
\item \wolfram|ProductFinalStates[A1_, A2_]| restituisce l'automa prodotto a partire da due automi (composizione parallela).
\item \wolfram|TuringMachineGraph[mu_]| restituisce il grafo della macchina di Turing.
\end{itemize}
\end{defi}
\begin{rema}
Nelle firme delle funzioni il carattere \wolfram{_} indica un \emph{pattern}, cioè un parametro formale. Quando si chiama la funzione si scrive solo l'argomento, senza il trattino basso.
\end{rema}
\subsection{Automi}
\begin{defi}
Un automa si specifica con tre oggetti:
\begin{itemize}
\item la lista \wolfram{T} delle transizioni, ciascuna della forma \wolfram|{stato, simbolo} -> nuovo_stato|;
\item lo stato iniziale \wolfram{q0};
\item l'insieme \wolfram{F} degli stati finali.
\end{itemize}
L'automa è la terna \wolfram|{T, q0, F}|.
\end{defi}
\begin{rema}
Esempio di automa con due stati, \wolfram{pari} e \wolfram{dispari}, sull'alfabeto $\{0,1\}$. Lo stato cambia a ogni simbolo $1$ letto e resta invariato a ogni simbolo $0$, quindi l'automa accetta le parole con un numero pari di $1$.
\begin{walfram}
T = {
  {pari, 0} -> pari,
  {dispari, 0} -> dispari,
  {pari, 1} -> dispari,
  {dispari, 1} -> pari
  }
q0 = pari
F = {pari}
automa = {T, q0, F}
\end{walfram}
Per mostrare l'esecuzione:
\begin{walfram}
init = InitAutomaton[
    automa2, {0, 1, 1, 0, 0, 1, 1, 1, 1, 1, 0, 0, 0, 1}
    ] ;
mu = DFA2TM[automa2] ;
TuringMachineGraph[mu]
ManipulateComputation[mu, init]
\end{walfram}
\end{rema}
\begin{rema}
Una volta definito \wolfram{automa}, lo si può inizializzare su una parola $w$ con \wolfram|InitAutomaton[automa, w]|, che restituisce la configurazione iniziale. Con \wolfram|DFA2TM[automa]| si ottiene la macchina di Turing equivalente, di cui si può visualizzare la computazione con \wolfram|ManipulateComputation| e il grafo con \wolfram|TuringMachineGraph|.
\end{rema}
\begin{rema}
\textbf{Composizione di automi:}
Per studiare sistemi complessi o l'intersezione di linguaggi regolari, la libreria mette a disposizione il prodotto cartesiano tramite le funzioni \wolfram{ProductTable[A1, A2]} e \wolfram{ProductFinalStates[A1, A2]}. 
La prima genera le transizioni incrociate combinando gli stati con il comando \wolfram{Tuples}, mentre la seconda definisce gli stati finali dell'automa prodotto.
\end{rema}

\begin{defi}
La funzione di visualizzazione interattiva
\[
\wolfram{ManipulateComputation[mu, init]}
\]
combina il motore di animazione di Wolfram (\wolfram{Manipulate}) con la generazione sequenziale delle configurazioni della macchina di Turing, permettendo di scorrere passo-passo l'intera computazione fino al limite di passi predefinito (\wolfram{maxsteps}).
\end{defi}